% Generated draft: requires mathematical and source review.
\documentclass[11pt]{article}
\usepackage[margin=25mm]{geometry}
\usepackage{fontspec}
\setmainfont{Latin Modern Roman}
\usepackage{amsmath,amssymb,mathtools,unicode-math}
\setmathfont{Latin Modern Math}
\usepackage{enumitem,needspace,xurl,hyperref}
\setlength{\parindent}{0pt}
\setlength{\parskip}{5pt}
\setlength{\emergencystretch}{3em}
\newcommand{\N}{\mathbb{N}}
\newcommand{\Z}{\mathbb{Z}}
\newcommand{\Q}{\mathbb{Q}}
\newcommand{\R}{\mathbb{R}}
\newcommand{\C}{\mathbb{C}}
\newcommand{\F}{\mathbb{F}}
\newcommand{\A}{\mathbb{A}}
\newcommand{\PP}{\mathbb{P}}
\newcommand{\E}{\mathbb{E}}
\newcommand{\eps}{\varepsilon}
\newcommand{\dd}{\,\mathrm{d}}
\begin{document}
\section*{UnsolvedMath ID: 11000266}
\textbf{Source problem code:} catalogue entry
\subsection{Definitions and mathematical statement}
Let \(S_g\) be the source's genus-\(g\) surface and \(\mathcal T(S_g)\) its Teichmüller space. Question 1.1 asks whether there is an equivariant deformation retraction onto a cocompact spine of dimension exactly
\[
 4g-5,
\]
identified in the source as \(\operatorname{vcd}(\operatorname{Mod}^{\pm}(S_g))\). The usual closed-surface formula uses \(g\ge2\); the main extract does not explicitly state the surface or genus hypotheses. The equivariance group should be confirmed, with \(\operatorname{Mod}^{\pm}(S_g)\) the natural extended-group reading of the notation.

In a usual strong-deformation-retract convention, this asks for an invariant spine \(K\subset\mathcal T(S_g)\) and a homotopy \(H_t\) with
\[
 H_0=\operatorname{id},\quad H_1(\mathcal T(S_g))=K,
 \quad H_t|_K=\operatorname{id}_K,\quad
 H_t(\gamma x)=\gamma H_t(x),
\]
with \(K/\operatorname{Mod}^{\pm}(S_g)\) compact and \(\dim K=4g-5\), subject to the source's spine/dimension conventions. An abstract contractible model of the same dimension, or a spine with noncompact quotient, is not automatically the requested equivariant retraction inside Teichmüller space.

The following context compares mapping class groups, \(\mathrm{GL}(n,\Z)\) and \(\operatorname{Out}(F_n)\), and points to Section 2.1 for related questions. It defines the Tits alternative by requiring every subgroup either to contain a nonabelian free group or to be virtually solvable. It reports proofs for all three families using ping-pong ideas on boundaries, with particularly intricate \(\operatorname{Out}(F_n)\) references [4,3] versus [9], while stressing that detailed techniques differ despite common proof architecture.

The source then describes a nonpositively curved invariant metric on the symmetric-space model for \(\mathrm{GL}(n,\Z)\), and two natural Teichmüller-space metrics, the Teichmüller metric and the one printed ``Weyl--Petersen'' (the usual Weil--Petersson terminology needs confirmation). It says Outer-space metric theory was undeveloped with no obvious natural metric candidate, a historical assessment rather than a current claim, and introduces a following vague question not included here. The isolated footnote about negative curvature for \(g\ge2\) has a missing antecedent and is not inserted as a curvature hypothesis on the requested spine.
\subsection{Short English statement}
Seek an equivariant, cocompact Teichmüller-space spine of dimension \(4g-5\), reached by deformation retraction. The source's comparisons with arithmetic groups and Outer space are retained as context, without replacing the embedded-spine request by an abstract action-dimension question.
\subsection{Sources}
\begin{itemize}
\item Local UnsolvedMath registry, permanent ID \(11000266\), exact source code \texttt{catalogue entry}, Question 1.1, PDF page 327. No individual proposer or proposed year is supplied; the title is truncated.
\item Farb, edited volume \emph{Problems on Mapping Class Groups and Related Topics}: \url{https://www.math.uchicago.edu/~farb/papers/mcgbook.pdf}. The volume has 329 explicitly typeset numbered items across chapters, repeated local numbering distinguished by PDF page. Extracted furniture identifies the automorphism-groups chapter and page 321, with an isolated footnote 1.
\item Internal references [4,3] and [9] for the compared Tits-alternative proofs, and cross-reference Section 2.1; full bibliography entries are absent.
\item Extraction pdf-font-structure, accessed 2026-07-29. Source evidence NEEDS\_REVIEW presents an open problem and requires a current primary-literature audit. Supplied OPEN-TRIAGE summary mentions Teichmüller actions and low-dimensional spines, relates them to action dimension, and explicitly cannot verify the exact question.
\item Dataset attribution: Ulam AI / UnsolvedMath Contributors, CC BY 4.0. Source rights NEEDS\_REVIEW: publicly accessible author/publisher copy, redistribution terms requiring review.
\end{itemize}
\end{document}
